\begin{lemma}[Values and bijectivity of a three-edge switch]
Let $f:A\to B$ and let $x,a,b\in A$ be pairwise distinct. For the map
$g=S(f;x,a,b)$ of \noderef{MatchingThreeEdgeSwitch},
\[
g(x)=f(a),\qquad g(a)=f(b),\qquad g(b)=f(x).
\]
For every $z\notin\{x,a,b\}$, $g(z)=f(z)$. If $f$ is bijective, then $g$
is bijective.
\end{lemma}
\begin{proof}
By \noderef{MatchingThreeEdgeSwitch}, first interchange $a,b$, then
interchange $x,a$, and finally apply $f$. The two interchanges send
$x$ to $a$, $a$ to $b$, and $b$ to $x$, since the three elements are
distinct. They fix every other element. This gives all four value
identities. Each interchange is its own inverse, so both are
bijections. Their composition with the bijection $f$ is a bijection.
\end{proof}
